\input{preamble}
\usepackage{eurosym} %les euros
\newcommand{\C}{\mathcal{C}}

\begin{document}

\pagestyle{fancy}
\fancyhead[L]{Seconde}
\fancyhead[C]{\textbf{TD n°5 : évolution chiffrée}}
\fancyhead[R]{\today}

\exe{}{
Dans chaque cas, dire si la quantité augmente, diminue ou reste constante.
\begin{enumerate}
\item La quantité est multipliée par un coefficient multiplicateur égal à $\frac13$.
\item La quantité est multipliée par un coefficient multiplicateur égal à $\frac14 + \frac13$.
\item La quantité subit deux évolutions successives de coefficient multiplicateur égal à $\frac12$.
\item La quantité est multipliée par un coefficient multiplicateur égal à $1,0001$.
\end{enumerate}
}{exe:1}{
Voir cours, il suffit de comparer le coefficient multiplicateur à 1.
}

\exe{}{
Dans chacun des cas ci-dessous, proposer des valeurs cohérentes avec les contraintes.
\begin{enumerate}
\item Les valeurs initiales et finales d'une quantité dont la variation absolue est de 42,5.
\item Les valeurs initiales et finales d'une quantité dont la variation relative est de 0,22.
\item La valeur initiale d'une quantité dont la variation absolue est de $-\frac{18}3$ et la valeur finale 5.
\item La valeur initiale d'une quantité dont la variation relative est de $\frac65$ et la valeur finale $8$.
\end{enumerate}
}{exe:2}{
\begin{enumerate}
\item On propose $x_1 = 0$ et $x_2=42,5$.
\item On propose $x_1 = 1$ et $x_2=1,22$.
\item On doit résoudre $5 - x_1 = -\frac{18}3$, on trouve $x_1=11$.
\item On doit résoudre $\frac{8-x_1}{x_1} = \frac65$ soit $8-x_1 = \frac65 x_1$. On trouve $x_1= \frac{40}{11}$.
\end{enumerate}
}

\exe{}{
	Calculer sans calculatrice les valeurs suivantes.
	\begin{multicols}{3}
	\begin{enumerate}
		\item $75\%$ de $60$
		\item $60\%$ de $75$
		\item $72\%$ de $25$
		\item $68\%$ de $20$
		\item $125\%$ de $40$
		\item $40\%$ de $125$
	\end{enumerate}
	\end{multicols}
}{exe:3}{
	\begin{multicols}{2}
	\begin{enumerate}
		\item $\dfrac34 \cdot 60 = 3 \cdot \dfrac{60}4 = 3 \cdot 15 = 45$
		\item $45$
		\item $\dfrac14 \cdot 72 = 18$
		\item $\dfrac15 \cdot 68 = \dfrac{136}{10} = 13,6$
		\item $40 + \dfrac14 \cdot 40 = 50$
		\item $50$
	\end{enumerate}
	\end{multicols}
}

\exe{}{
Une quantité subit deux évolutions successives de coefficients multiplicateurs $m_1$ et $m_2$.
\begin{enumerate}
\item Donner la formule permettant de calculer le coefficient multiplicateur total.
\item Donner la formule permettant de calculer une proportion d'évolution à partir d'un coefficient multiplicateur.
\item Rappeler comment convertir une proportion en pourcentage.
\item Application numérique : calculer le pourcentage d'évolution d'une quantité subissant deux évolutions successives 
de coefficients multiplicateurs $m_1 = 1,5$ et $m_2 = 1,2$.
\end{enumerate}
}{exe:4}{
\begin{enumerate}
\item $m = m_1 \times m_2$.
\item$p = m-1$.
\item $p \% = p \times 100$.
\item Application numérique : $m = 1,5 \times 1,2 = 1,8$ donc $p = 0,8$ et $p \% = 80 \%$. 
Ce qui correspond à une augmentation de $80 \%$.
\end{enumerate}
}

\exe{}{
Dans chacun des cas ci-dessous, donner le taux d'évolution total.
\begin{enumerate}[label=(\alph*)]
\item Une augmentation de $20 \%$ suivie d'une augmentation de $15 \%$.
\item Une augmentation de $10 \%$ suivie d'une baisse de $10 \%$.
\item Une baisse de $20 \%$ suivie d'une augmentation de $25 \%$.
\item Une augmentation d'une proportion 0,25 suivie d'une baisse d'une proportion 0,20.
\end{enumerate}
}{exe:5}{
\begin{enumerate}[label=(\alph*)]
\item On transforme d'abord chaque pourcentage d'évolution en coefficient multiplicateur.
Avec les formules de l'exercice $\ref{exe:4}$ on trouve :
\[ m_1 = 1 + \frac{p_1 \%}{100} = 1,2 \quad \text{et} \quad m_2 = 1 + \frac{p_2 \%}{100} = 1,15 . \]
Le coefficient multiplicateur correspond est donc $m = 1,2 \times 1,15 = 1,38$. 
Ce qui correspond à une évolution d'une proportion $p = m - 1 = 0,38$,
soit une augmentation de $38 \%$.
\item On procède de la même façon en prenant garde que, pour une baisse, la proportion est négative.
Avec les formules de l'exercice $\ref{exe:4}$ on trouve :
\[ m_1 = 1 + \frac{p_1 \%}{100} = 1,1 \quad \text{et} \quad m_2 = 1 + \frac{p_2 \%}{100} = 1 - 0,1 = 0,9 . \]
Donc $m = 0,99$, soit $p=-0,01$ ce qui correspond à une baisse de $1 \%$.
\item On procède de même et on trouve $m=1$ soit $p=0$, pas d'évolution.
\item Même résultat.
\end{enumerate}
}

\exe{}{
Dans chacun des cas ci-dessous, donner le taux d'évolution réciproque.
\begin{enumerate}[label=(\alph*)]
\item Une augmentation de $20 \%$.
\item Une baisse d'une proportion $0,7$.
\item Une augmentation d'une proportion $0,45$.
\item Une augmentation de $33 \%$.
\end{enumerate}
}{exe:6}{
\begin{enumerate}[label=(\alph*)]
\item Une augmentation de $20 \%$ correspond à un coefficient multiplicateur de 1,2.
Le coefficient multiplicateur réciproque est donc égal à $m' = \frac{1}{1,2} = \frac{10}{12} = \frac{5}{6}$.
Ceci correspond à une proportion d'évolution de $p' = m' - 1 = - \frac16$ soit une \textbf{baisse} d'environ $16,7 \%$.
\item On emploie la même méthode, on trouve une augmentation d'environ $233 \%$.
\item Une baisse d'environ $31 \%$.
\item Une baisse d'environ $25 \%$.
\end{enumerate}
}

\exe{}{
  En 2023 en France, $13\%$ des espèces (faune et flore) sont considérées comme menacées à l'échelle mondiale (catégories ``vulnérable'' à ``danger critique'' de l'UICN).
  Parmis celles-ci, $23\%$ sont en danger critique.
  
  Calculer le pourcentage d'espèces en danger critique par rapport au nombre total d'espèces.
}{exe:7}{
	Notons $E$ l'ensemble des espèces indigènes à la France, $M$ la sous-population des espèces menacées, et $C$ la sous-population des espèces en danger critique.
	On a donc la suite d'inclusions
		\[ C \subset M \subset E. \]
	
	Le texte donne les informations
		\begin{align*}
			p_M = 0,13 && \text{et} && p_{C / M }= 0,23.
		\end{align*}
	Selon la formule du cours, $p_C = p_M \times p_{C / M} = 0,23 \times 0,13 = 0,0299$.
	Donc les espèces en danger critique constituent $2,99 \%$ des espèces.

	Les proportions sont ainsi multiplicatives. Attention à ne pas naïvement multiplier les pourcentages, car $13\times 23 = 299$.
}

\exe{, difficulty=1}{
	Lucille dépose 100 \euro~ à la banque sur un compte rémunéré à $3 \%$ par an. 
	Ce qui signifie que, chaque année, elle recevra un versement correspondant à $3 \%$ de la somme présente sur le compte.
	
	\begin{enumerate}
		\item Combien d'argent aura-t-elle après $1$ ans ?
		\item Combien d'argent aura-t-elle après $2$ ans ?
		\item Combien d'argent y aura-t-il sur son compte après $100$ ans ?
	\end{enumerate}
	
	%\hfill
}{exe:8}{
	\begin{enumerate}
		\item On multiplie $100$ par $1,03$, ce qui donne $100 \times 1,03 = 103$ \euro.
		\item On multiplie $100$ par $1,03$ puis on multiplie ce résultat par $1,03$, ce qui revient à calculer $100 \times 1,03^2$.
		\item En suivant le même raisonnement, on multiplie $100$ fois par $1,03$, ce qui donne 
		$100 \times 1,03^{100} \approx 1921,86$ \euro.
	\end{enumerate}
	Lucille a multiplié son placement par plus de 19 !
}

\exe{, difficulty=1}{
Un commerçant diminue le prix d'un article de 1~000 \euro~à 999 \euro. 
Calculer le pourcentage d'évolution nécessaire pour que le prix de l'article revienne à sa valeur initiale.
}{exe:9}{
Le coefficient multiplicateur correspond à la baisse du prix est $m = \frac{999}{1000}$. 
Le coefficient multiplicateur réciproque est donc égal à $m' = \frac{1}{m} = \frac{1000}{999}$.
Ceci correspond à une proportion de $p' = m' - 1 = \frac1{999}$ soit une augmentation d'environ $0,1001 \%$.
}

\exe{, difficulty=2}{
	Une livre anglaise équivaut à 0,453 kilogrammes environ.
	Pour convertir une quantité $x$ de livres anglaises en kilogrammes, la procédure suivante est proposée.
		\begin{enumerate}[label=\roman*)]
			\item Prendre d'abord la moitié de $x$ ; puis
			\item Diminuer le résultat de 10\%.
		\end{enumerate}
	\begin{enumerate}
		\item 
		Convertir 10 livres anglaises en kilogrammes de deux façon différentes,
			\begin{enumerate}[label=(\alph*)]
				\item en utilisant que 1 livre = 0,453 kilogrammes ; et
				\item en utilisant la procédure proposée.
			\end{enumerate}
		\item 
		Calculer le coefficient multiplicateur qui permet de convertir une quantité de livres anglaises en kilogrammes
			\begin{enumerate}[label=(\alph*)]
				\item lorsqu'on utilise que 1 livre = 0,453 kilogrammes ; et
				\item lorsqu'on utilise la procédure proposée.
			\end{enumerate}
		\item
		Donner les évolutions associées aux deux coefficient multiplicateurs.
		\item
		Calculer l'erreur relative de la procédure proposée : quelle évolution permet d'obtenir la valeur de la procédure à partir de la valeur réelle de conversion ?
		\item
		Montrer que l'erreur absolue, elle, peut être aussi grande que voulue.
		L'erreur absolue entre deux valeurs est la différence entre la plus grande et la plus petite valeur.
	\end{enumerate}
}{exe:10}{
	\begin{enumerate}
		\item 
			\begin{enumerate}[label=(\alph*)]
				\item
				Par proportionnalité, 10 livres équivalent à 4,53 kilogrammes.
				\item 
				En utilisant la procédure proposée, on diviser 10 par 2 pour trouver 5, puis on soustrait 10\%, ce qui donne $5-0,5 = 4,5$ kilogrammes.
			\end{enumerate}
		\item 
		Calculer le coefficient multiplicateur qui permet de convertir une quantité de livres anglaises en kilogrammes
			\begin{enumerate}[label=(\alph*)]
				\item 
				Par proportionnalité, $x$ livres équivaut à $0,453x$ kilogrammes : le coefficient est $0,453$.
				\item
				La procédure calcule d'abord $0,5x$ puis $0,9 \times 0,5 x = (0,9 \times 0,5) x = 0,45x$.
				Le coefficient est donc $0,45$.
			\end{enumerate}
		\item
		Multiplier par $0,453$ équivaut à une diminution de $54,7\%$.
		Multiplier par $0,45$ équivaut à une diminution de $55\%$.
		\item
		On cherche le coefficient $CM$ qui vérifie $CM \times 0,453x = 0,45x \iff CM = \frac{0,45}{0,453} \approx 0,993$
		Ce coefficient correspond à une diminution de $0,7\%$ environ : c'est l'erreur relative commise par la procédure.
		\item
		On calcule ici $0,453x - 0,45x = 0,003x$, qui peut être aussi grand que voulu, pourvu que $x$ soit très grand également.
		Par exemple, pour obtenir une erreur de $10~000$, il faut prendre $x=\frac{10~000}{0,003} \approx 3~333~334$.
		
		L'erreur relative, quant à elle, sera toujours de $0,7\%$ environ.
		Cependant, $0,7\%$ d'une très grande valeur est toujours une très grande valeur !
	\end{enumerate}
}

%%%%%%%%%%%%

\newpage
\fancyhead[C]{\textbf{Solutions}}
\shipoutAnswer



\end{document}
